% Part: second-order-logic
% Chapter: sol-and-set-theory
% Section: comparing-sets

\documentclass[../../../include/open-logic-section]{subfiles}

\begin{document}

\olfileid{sol}{set}{cmp}

\olsection{مقارنة المجموعات}

\begin{prop}
تعرّف الـ!!{formula}~$\lforall[x][(X(x) \lif Y(x))]$ علاقة الاحتواء
الجزئي؛ أي إن~$\Sat{M}{\lforall[x][(X(x) \lif Y(x))]}[s]$ إذا وفقط
إذا كان~$s(X) \subseteq s(Y)$.
\end{prop}

\begin{prop}
تعرّف الـ!!{formula}~$\lforall[x][(X(x) \liff Y(x))]$ علاقة الهوية
على المجموعات؛ أي إن~$\Sat{M}{\lforall[x][(X(x) \liff Y(x))]}[s]$
إذا وفقط إذا كان~$s(X) = s(Y)$.
\end{prop}

\begin{prop}
تعرّف الـ!!{formula}~$\lexists[x][X(x)]$ خاصية عدم الخلو؛ أي إن
$\Sat{M}{\lexists[x][X(x)]}[s]$ إذا وفقط إذا كان~$s(X) \neq
\emptyset$.
\end{prop}

لا تكون المجموعة~$X$ أكبر من المجموعة~$Y$،
$\cardle{X}{Y}$، إذا وفقط إذا وجدت دالة~!!a{injective}
$f\colon X \to Y$. ولما كان يمكننا التعبير عن كون دالة ما متباينة، وعن
كون قيمها عند وسائط من~$X$ تنتمي إلى~$Y$ أيضًا، أمكننا كذلك تعريف
علاقة عدم الكبر على المجموعات الجزئية من المجال.

\begin{prop}
تعرّف الصيغة
\[
\lexists[u][(\lforall[x][(X(x) \lif Y(u(x)))] \land \lforall[x][\lforall[y][(\eq[u(x)][u(y)] \lif
      \eq[x][y])]])]
\]
علاقة عدم الكبر.
\end{prop}

تكون مجموعتان متساويتين في الحجم، أو «متكافئتين عدديًا»،
$\cardeq{X}{Y}$، إذا وفقط إذا وجدت دالة~!!a{bijective}
$f\colon X \to Y$.

\begin{prop}
تعرّف الصيغة
\begin{multline*}
  \lexists[u][(\lforall[x][(X(x) \lif Y(u(x)))] \land {}\\
    \lforall[x][\lforall[y][((X(x) \land X(y) \land \eq[u(x)][u(y)]) \lif \eq[x][y])]]
  \land {} \\\lforall[y][(Y(y) \lif \lexists[x][(X(x)
      \land \eq[y][u(x)])])])]
\end{multline*}
علاقة التكافؤ العددي.
\end{prop}

وسنختصر هذه الـ!!{formula}s، على الترتيب، بالرموز
$X \subseteq Y$، و$X = Y$، و$X \neq \emptyset$، و$\cardle{X}{Y}$،
و$\cardeq{X}{Y}$. (وقد يكون هذا مربكًا قليلًا، لأننا نستعمل الترميز
نفسه حين نتحدث بصورة غير صورية عن المجموعتين~$X$ و$Y$؛ أما هنا فالترميز
اختصار لـ!!{formula}s من منطق الرتبة الثانية تتضمن متغيري العلاقة
الأحاديي الموضع~$X$ و$Y$.)

\begin{prop}
الـ!!{sentence}~$\lforall[X][\lforall[Y][((\cardle{X}{Y} \land
    \cardle{Y}{X}) \lif \cardeq{X}{Y})]]$ صادقة منطقيًا.
\end{prop}

\begin{proof}
تُرضى الـ!!{sentence} في~!!a{structure}~$\Struct{M}$ إذا كان، لكل
مجموعتين جزئيتين~$X \subseteq \Domain{M}$ و
$Y \subseteq \Domain{M}$، أن~$\cardeq{X}{Y}$ متى كان
$\cardle{X}{Y}$ و$\cardle{Y}{X}$. لكن هذا يصدق على \emph{أي}
مجموعتين~$X$ و$Y$؛ فهو مبرهنة شرودر--برنشتاين.
\end{proof}

\end{document}
